\documentclass[11pt,a4paper]{article}

\usepackage[margin=0.72in]{geometry}
% Libertinus from the source template is not installed in this TeX Live setup.
% Latin Modern is available and keeps the document compilable here.
\usepackage{lmodern}
\usepackage{microtype}
\usepackage{xcolor}
\usepackage{enumitem}
\usepackage{tabularx}
\usepackage{array}
\usepackage{booktabs}
\usepackage{longtable}
\usepackage{amsmath}
\usepackage{amssymb}
\usepackage{fancyhdr}
\usepackage{titlesec}
\usepackage{hyperref}
\usepackage{multicol}
\usepackage{parskip}

\definecolor{navy}{HTML}{183153}
\definecolor{blue}{HTML}{2563EB}
\definecolor{green}{HTML}{166534}
\definecolor{red}{HTML}{B91C1C}
\definecolor{gold}{HTML}{92400E}
\definecolor{lightblue}{HTML}{EFF6FF}
\definecolor{lightgreen}{HTML}{F0FDF4}
\definecolor{lightred}{HTML}{FEF2F2}
\definecolor{lightgold}{HTML}{FFFBEB}

\hypersetup{colorlinks=true,linkcolor=blue,urlcolor=blue}
\pagestyle{fancy}
\fancyhf{}
\lhead{Probability, Statistics and Computational Mathematics}
\rhead{Modules 4--5 Exam Notes}
\cfoot{\thepage}

\titleformat{\section}{\Large\bfseries\color{navy}}{}{0pt}{}
\titleformat{\subsection}{\large\bfseries\color{navy}}{}{0pt}{}
\titleformat{\subsubsection}{\normalsize\bfseries\color{blue}}{}{0pt}{}
\setlist[itemize]{leftmargin=*,nosep}
\setlist[enumerate]{leftmargin=*,nosep}

\newenvironment{exambox}{\par\medskip\noindent\color{blue}\textbf{Exam Focus}\par\smallskip\color{black}\begin{quote}}{\end{quote}\medskip}
\newenvironment{rememberbox}{\par\medskip\noindent\color{gold}\textbf{Remember}\par\smallskip\color{black}\begin{quote}}{\end{quote}\medskip}
\newenvironment{dangerbox}{\par\medskip\noindent\color{red}\textbf{Do Not Confuse}\par\smallskip\color{black}\begin{quote}}{\end{quote}\medskip}
\newenvironment{answerbox}{\par\medskip\noindent\color{green}\textbf{Answer Skeleton}\par\smallskip\color{black}\begin{quote}}{\end{quote}\medskip}
\newcommand{\mkmarks}[1]{\hfill\textbf{[#1 marks]}}

\begin{document}

\begin{center}
{\Huge\bfseries\color{navy} PROBABILITY, STATISTICS AND COMPUTATIONAL MATHEMATICS}\\[4pt]
{\LARGE\bfseries MODULES 4 AND 5}\\[8pt]
{\large High-yield notes, formulas, worked examples, and question bank}\\[10pt]
\rule{0.85\textwidth}{0.6pt}
\end{center}

\begin{exambox}
\textbf{Scope.} These notes cover only Module 4 (Mathematical Logic) and Module 5 (Counting and Mathematical Induction), matching the supplied module PDFs. The exam timetable lists the MCA101 paper for Friday, 30 October 2026, 9:30--10:30 a.m.; the marks pattern was not supplied, so the question lengths below are revision guidance rather than a confirmed blueprint. For proof and calculation questions, show the rule or formula, substitute values, and state the conclusion clearly.
\end{exambox}

\tableofcontents
\newpage

\section{Module 4: Mathematical Logic}

\subsection{4.1 Propositions and connectives}
A \textbf{proposition} is a declarative statement with exactly one truth value: true (T) or false (F). Questions, commands, and open sentences containing an unspecified variable are not propositions.

\begin{center}\small
\begin{tabularx}{\textwidth}{>{\bfseries}p{0.16\textwidth} p{0.18\textwidth} X X}
\toprule
Connective & Notation & Meaning & Truth condition\\
\midrule
Negation & $\neg p$ & not $p$ & True exactly when $p$ is false\\
Conjunction & $p\land q$ & $p$ and $q$ & True when both are true\\
Inclusive OR & $p\lor q$ & $p$ or $q$ or both & False only when both are false\\
Exclusive OR & $p\oplus q$ & exactly one of $p,q$ & True when their values differ\\
Implication & $p\to q$ & if $p$, then $q$ & False only for $p=T,q=F$\\
Biconditional & $p\leftrightarrow q$ & $p$ iff $q$ & True when their values agree\\
\bottomrule
\end{tabularx}\end{center}

For $n$ variables, a complete truth table has $2^n$ rows. The implication $p\to q$ is logically equivalent to $\neg p\lor q$, and $p\leftrightarrow q$ is equivalent to $(p\to q)\land(q\to p)$.

\begin{center}\small
\begin{tabular}{cc|cccccc}
\toprule $p$&$q$&$\neg p$&$p\land q$&$p\lor q$&$p\oplus q$&$p\to q$&$p\leftrightarrow q$\\\midrule
T&T&F&T&T&F&T&T\\
T&F&F&F&T&T&F&F\\
F&T&T&F&T&T&T&F\\
F&F&T&F&F&F&T&T\\\bottomrule
\end{tabular}\end{center}

\subsection{4.2 Tautology, contradiction, contingency}
\begin{itemize}
\item A \textbf{tautology} is true on every row (example: $p\lor\neg p$).
\item A \textbf{contradiction} is false on every row (example: $p\land\neg p$).
\item A \textbf{contingency} is true on some rows and false on others (example: $p\land q$).
\end{itemize}
To classify a formula, construct its full truth table and inspect the final column.

\subsection{4.3 Logical equivalences}
Two formulas are logically equivalent, written $P\equiv Q$, when their truth values agree for every assignment. Use a truth table or the following laws.

\begin{center}\small
\begin{tabularx}{\textwidth}{>{\bfseries}p{0.22\textwidth} X X}
\toprule Law & Form 1 & Form 2\\\midrule
Identity & $p\land T$ & $p$; \quad $p\lor F=p$\\
Domination & $p\lor T$ & $T$; \quad $p\land F=F$\\
Idempotent & $p\lor p$ & $p$; \quad $p\land p=p$\\
Double negation & $\neg\neg p$ & $p$\\
Commutative & $p\lor q$ & $q\lor p$; similarly for $\land$\\
Associative & $(p\lor q)\lor r$ & $p\lor(q\lor r)$; similarly for $\land$\\
Distributive & $p\land(q\lor r)$ & $(p\land q)\lor(p\land r)$\\
De Morgan & $\neg(p\land q)$ & $\neg p\lor\neg q$\\
De Morgan & $\neg(p\lor q)$ & $\neg p\land\neg q$\\
Absorption & $p\lor(p\land q)$ & $p$; dual: $p\land(p\lor q)=p$\\
Implication & $p\to q$ & $\neg p\lor q$\\
Biconditional & $p\leftrightarrow q$ & $(p\to q)\land(q\to p)$\\
\bottomrule
\end{tabularx}\end{center}

\subsubsection{Converse, inverse, contrapositive}
For $p\to q$: the \textbf{converse} is $q\to p$; the \textbf{inverse} is $\neg p\to\neg q$; the \textbf{contrapositive} is $\neg q\to\neg p$. The original and contrapositive are equivalent. The converse and inverse are equivalent to each other, but are not generally equivalent to the original.

\begin{dangerbox}For “if $p$, then $q$,” do not swap the hypothesis and conclusion and assume the result is equivalent. Only the contrapositive is guaranteed equivalent.\end{dangerbox}

\subsection{4.4 Predicates and quantifiers}
A \textbf{predicate} $P(x)$ is a statement whose truth depends on the value assigned to $x$. A domain specifies the allowed values. Quantifiers turn predicates into propositions:
\begin{itemize}
\item $\forall x\,P(x)$: for every $x$ in the domain, $P(x)$ is true.
\item $\exists x\,P(x)$: there is at least one $x$ in the domain for which $P(x)$ is true.
\end{itemize}
Negation rules: $\neg\forall x\,P(x)\equiv\exists x\,\neg P(x)$ and $\neg\exists x\,P(x)\equiv\forall x\,\neg P(x)$. For example, “not everyone passed” means “someone did not pass,” not “nobody passed.”

\subsubsection{Nested quantifiers}
Quantifier order can change meaning. Over a suitable domain, $\forall x\exists y\,R(x,y)$ says each $x$ may have its own $y$; $\exists y\forall x\,R(x,y)$ says one fixed $y$ works for every $x$. In general, these are not equivalent. To disprove a universal statement, give a counterexample; to prove an existential statement, give a witness.

\subsection{4.5 Rules of inference}
An argument is valid when its conclusion follows from its premises. Common valid forms:
\begin{center}\small
\begin{tabularx}{\textwidth}{>{\bfseries}p{0.24\textwidth} X X}
\toprule Rule & Premises & Conclusion\\\midrule
Modus ponens & $p\to q$, $p$ & $q$\\
Modus tollens & $p\to q$, $\neg q$ & $\neg p$\\
Hypothetical syllogism & $p\to q$, $q\to r$ & $p\to r$\\
Disjunctive syllogism & $p\lor q$, $\neg p$ & $q$\\
Addition & $p$ & $p\lor q$\\
Simplification & $p\land q$ & $p$ (or $q$)\\
Conjunction & $p$, $q$ & $p\land q$\\
Resolution & $p\lor q$, $\neg p\lor r$ & $q\lor r$\\
\bottomrule
\end{tabularx}\end{center}

\subsection{4.6 Normal forms}
An \textbf{elementary product} (minterm) is an AND of literals; an \textbf{elementary sum} (maxterm) is an OR of literals. A literal is a variable or its negation.
\begin{itemize}
\item \textbf{DNF}: OR of AND terms, e.g. $(p\land q)\lor(\neg p\land r)$.
\item \textbf{CNF}: AND of OR clauses, e.g. $(p\lor q)\land(\neg p\lor r)$.
\end{itemize}
To convert algebraically, first remove $\to$ and $\leftrightarrow$, push negations inward using De Morgan's laws, then distribute. For canonical forms, use truth-table rows: DNF takes the rows where the formula is true (minterms); CNF takes rows where it is false (maxterms).

\begin{exambox}\textbf{Worked conversion.} $p\to q\equiv\neg p\lor q$, which is already CNF. To write it as DNF, expand with $x\equiv(x\land T)$ and distribute, or use the true rows: $(\neg p\land\neg q)\lor(\neg p\land q)\lor(p\land q)$.\end{exambox}

\subsection{4.7 Exam answer patterns}
\begin{answerbox}
\textbf{Truth table / tautology:} list variables; make all $2^n$ assignments; add intermediate connective columns; evaluate the target column; classify and state the result.\\
\textbf{Equivalence proof:} start with one side; name each identity used; transform step by step to the other side.\\
\textbf{Inference:} number premises; identify the rule applied at each line; show the conclusion as the last line.\\
\textbf{Quantifier negation:} negate the quantifier, switch $\forall\leftrightarrow\exists$, then negate the predicate.
\end{answerbox}

\section{Module 5: Counting and Mathematical Induction}

\subsection{5.1 Addition and multiplication principles}
\textbf{Addition rule:} if a choice can be made in one of mutually exclusive cases with $a_1,a_2,\ldots,a_k$ outcomes, total outcomes are $a_1+a_2+\cdots+a_k$. If cases overlap, remove duplicate counts with inclusion--exclusion.

\textbf{Multiplication rule:} if a process has successive choices with $a_1,a_2,\ldots,a_k$ possibilities at each stage, total outcomes are $a_1a_2\cdots a_k$. The number may vary by stage; multiply the actual choices available at each step.

\begin{exambox}\textbf{Examples.} Three-letter strings followed by three digits, repetition allowed: $26^3\cdot10^3$. Select one boy from 28 and one girl from 18: $28\cdot18=504$. A die rolled five times: $6^5=7776$.\end{exambox}

\subsection{5.2 Factorials and permutations}
$n!=n(n-1)\cdots2\cdot1$, and $0!=1$. A permutation is an ordered arrangement; changing order creates a different result.
\[
{}^nP_r=P(n,r)=\frac{n!}{(n-r)!}\quad\text{(choose and arrange $r$ distinct objects from $n$).}
\]
If repetition is allowed and each of $r$ positions has $n$ choices, there are $n^r$ strings. If all $n$ objects are arranged with repeated indistinguishable types of multiplicities $n_1,\ldots,n_k$, the number is $\frac{n!}{n_1!\cdots n_k!}$.

\begin{exambox}\textbf{Worked examples.} Four distinct letters selected and ordered from six: $6P4=6!/2!=360$. Arrangements of ROOTS: $5!/2!=60$ (two O's). Six-digit PIN with digits 0--9 and repetition: $10^6$; leading zero is allowed for a PIN.\end{exambox}

\subsection{5.3 Combinations}
A combination is a selection where order does not matter:
\[
{n\choose r}=C(n,r)=\frac{n!}{r!(n-r)!},\qquad nP r={n\choose r}r!.
\]
Useful identities: $\binom n0=\binom nn=1$; $\binom n1=n$; $\binom nr=\binom n{n-r}$; $\binom nr+\binom n{r-1}=\binom{n+1}r$; $\sum_{r=0}^n\binom nr=2^n$.

\begin{exambox}\textbf{Worked examples.} Teams of 3 from 8: $\binom83=56$. A committee of 3 from 6 faculty including both named professors: choose the third member from the other four, so 4. From 5 boys and 4 girls, 3-person committees with at least one girl: $\binom93-\binom53=84-10=74$.\end{exambox}

\begin{dangerbox}Use permutations when order or role assignment matters; use combinations when only membership matters. A PIN permits leading zero; a number usually does not. “At least one” is often quickest by counting the total and subtracting the zero-case.\end{dangerbox}

\subsection{5.4 Pigeonhole principle}
\textbf{Basic principle:} placing $N$ objects into $k$ boxes with $N>k$ guarantees a box with at least two objects. \textbf{Generalized principle:} some box contains at least $\left\lceil N/k\right\rceil$ objects. To guarantee at least $r$ objects in one of $k$ boxes, the minimum required is
\[
k(r-1)+1.
\]
\textbf{Example:} 47 students in 8 groups gives a group of at least $\lceil47/8\rceil=6$. To guarantee 5 birthdays in the same month, $12(5-1)+1=49$ people suffice.

\subsection{5.5 Inclusion--exclusion principle}
For two finite sets:
\[
|A\cup B|=|A|+|B|-|A\cap B|.
\]
For three sets:
\[
|A\cup B\cup C|=|A|+|B|+|C|-|A\cap B|-|A\cap C|-|B\cap C|+|A\cap B\cap C|.
\]
For exactly two sets among three, add the three pairwise intersections and subtract three times the triple intersection. For exactly one, add the single-set totals, subtract twice the pairwise intersections, then add three times the triple intersection.

\begin{exambox}\textbf{Worked examples.} Of 80 students, 50 use Android and 45 iPhone, and everyone uses at least one: both $=50+45-80=15$. For the supplied language survey, at least one $=50+45+40-20-18-15+8=90$.\end{exambox}

\subsection{5.6 Mathematical induction}
To prove $P(n)$ for all integers $n\ge n_0$:
\begin{enumerate}
\item State $P(n)$ and the starting value $n_0$.
\item \textbf{Base case:} verify $P(n_0)$ directly.
\item \textbf{Inductive hypothesis:} assume $P(k)$ is true for an arbitrary $k\ge n_0$.
\item \textbf{Inductive step:} use that assumption to prove $P(k+1)$.
\item Conclude by induction that $P(n)$ holds for every integer $n\ge n_0$.
\end{enumerate}

\begin{exambox}\textbf{Worked proof.} Prove $1+2+\cdots+n=\frac{n(n+1)}2$ for $n\ge1$. Base: $n=1$ gives $1=1$. Assume $1+\cdots+k=\frac{k(k+1)}2$. Then $1+\cdots+k+(k+1)=\frac{k(k+1)}2+(k+1)=\frac{(k+1)(k+2)}2$. This is the formula for $k+1$; hence it holds for all $n\ge1$.\end{exambox}

\begin{dangerbox}In the inductive step, assume the claim only for $k$, then derive the $k+1$ case. Do not assume the statement you are trying to prove for $k+1$. Include both the base case and the final conclusion.\end{dangerbox}

\subsection{5.7 Quick method selector}
\begin{center}\small
\begin{tabularx}{\textwidth}{>{\bfseries}p{0.32\textwidth} X}
\toprule Clue in question & Method\\\midrule
“either ... or” in disjoint cases & Addition rule\\
several stages / “and then” & Multiplication rule\\
arrange, order, rank, positions & Permutation\\
choose, select, committee, team & Combination\\
guarantee a repeated category & Pigeonhole principle\\
overlapping groups / at least one set & Inclusion--exclusion\\
prove for every natural number & Mathematical induction\\
\bottomrule
\end{tabularx}\end{center}

\section{Important Exam Topics}
\begin{itemize}
\item Truth tables; tautology, contradiction, contingency; propositional equivalence laws.
\item Converse, inverse, contrapositive; translating statements; predicates, quantifiers, and quantifier negation.
\item Rules of inference and DNF/CNF conversion.
\item Addition and multiplication rules; factorials; permutations with/without repetition; combinations and committee constraints.
\item Generalized pigeonhole principle; two- and three-set inclusion--exclusion.
\item Mathematical induction proofs, especially sums and divisibility.
\end{itemize}

\section{Answer Templates}
\subsection{Prove a logical equivalence \mkmarks{5--10}}
\begin{answerbox}Write the expression; remove implications; apply the named equivalence law one step at a time; reach the target expression. Alternatively, show identical truth-table columns. Name each law to make the reasoning clear.\end{answerbox}

\subsection{Solve a counting problem \mkmarks{5--10}}
\begin{answerbox}Identify whether order matters and whether repetition is allowed; define the cases or stages; state the counting rule/formula; substitute values; simplify; answer in context. For restrictions, count valid cases only or use complement counting.\end{answerbox}

\subsection{Prove by induction \mkmarks{5--10}}
\begin{answerbox}Let $P(n)$ be the claim. Verify the base case. Assume $P(k)$. Starting from the left side for $k+1$, use the inductive hypothesis and simplify to the required right side. State that the claim follows for all specified $n$.\end{answerbox}

\section{Question Bank}
\subsection*{Definitions / short answers}
\begin{enumerate}
\item Define proposition, predicate, tautology, contradiction, and contingency.
\item State the truth condition for implication and biconditional.
\item State De Morgan's laws and the negation rules for quantifiers.
\item Define DNF and CNF, with one example of each.
\item Define permutation and combination; state the relation between them.
\item State the generalized pigeonhole principle and two-set inclusion--exclusion formula.
\end{enumerate}
\subsection*{Practice problems}
\begin{enumerate}
\item Construct truth tables for $p\lor\neg q$ and $(p\lor q)\lor\neg p$; classify each.
\item Use equivalence laws to show $p\to(q\to p)$ is a tautology.
\item Find the converse, inverse, and contrapositive of “If a number is divisible by 4, then it is even.”
\item Negate: “Every student solved some problem.” Express the result with quantifiers.
\item Identify the inference rule: $p\to q$, $q\to r$, therefore $p\to r$.
\item Convert $p\to(q\land r)$ to CNF and DNF.
\item How many 4-digit numbers can be made from digits 1--9 without repetition?
\item How many arrangements of INDEPENDENCE are there? How many start with P? (Use letter multiplicities carefully.)
\item From 6 men and 5 women, how many 4-person committees include at least one woman?
\item How many integers from 1 to 100 are divisible by 2 or 5?
\item What minimum number of people guarantees that at least 4 were born on the same day of the week?
\item In a class, 35 study Statistics, 28 study Discrete Mathematics, and 12 study both. Find at least one, only Statistics, only Discrete Mathematics, and exactly one.
\item Prove by induction that $1+3+5+\cdots+(2n-1)=n^2$ for $n\ge1$.
\item Prove by induction that $3\mid(n^3+2n)$ for every positive integer $n$.
\end{enumerate}

\section{Last-Minute Revision Page}
\begin{multicols}{2}
\textbf{LOGIC}\\
$p\to q\equiv\neg p\lor q$\\
$p\leftrightarrow q\equiv(p\to q)\land(q\to p)$\\
$\neg\forall xP\equiv\exists x\neg P$\\
$\neg\exists xP\equiv\forall x\neg P$\\[4pt]
\textbf{COUNTING}\\
Addition: sum disjoint cases\\
Multiplication: product successive choices\\
$nP r=n!/(n-r)!$\\
$nC r=n!/[r!(n-r)!]$\\
Repetition: $n^r$\\
Identical objects: $n!/\prod n_i!$\\[4pt]
\textbf{GUARANTEES / OVERLAP}\\
Pigeonhole: at least $\lceil N/k\rceil$\\
Guarantee $r$: $k(r-1)+1$\\
$|A\cup B|=|A|+|B|-|A\cap B|$\\
Three-set pattern: singles $-$ pairs $+$ triple\\[4pt]
\textbf{INDUCTION}\\
Base case $\to$ assume $P(k)$ $\to$ prove $P(k+1)$ $\to$ conclude\\[4pt]
\textbf{DO NOT CONFUSE}\\
Order: permutation; no order: combination.\\
Contrapositive equivalent to original; converse generally not.\\
At least one: consider complement.\\
\end{multicols}

\section*{Source Scope}
Prepared from the supplied \texttt{Module 4 Part 1.pdf} and \texttt{Module 5.pdf} in the Statistics folder, following the structure and visual style of \texttt{notes-template.tex}. Module 4's listed syllabus also includes predicates/quantifiers, nested quantifiers, and rules of inference; concise standard definitions and rules are included to cover those listed topics where the slides provide limited detail. Verify any instructor-specific conventions against class notes.

\vfill
\begin{center}\textit{Study principle: identify the rule, show the working, and state the conclusion.}\end{center}
\end{document}
